\newcommand{\corr}{\mathbbm{1}_{\text{correct}}}
\newcommand{\one}{\mathbf 1}
\newcommand{\ip}[1]{\left\langle #1 \right\rangle}
\newcommand{\EE}{\mathbb E}
\newcommand{\PP}{\mathbb P}
\newcommand{\sig}{\sigma}
\newcommand{\BR}{\mathrm{BR}}
\newcommand{\argmax}{\mathop{\mathrm{arg\,max}}}
\newcommand{\argmin}{\mathop{\mathrm{arg\,min}}}
\newcommand{\dH}{\pi_H}
\newcommand{\dS}{\pi_S}
\newcommand{\dHt}{\tilde{\pi}_H}
\newcommand{\dSt}{\tilde{\pi}_S}


\section*{1 \; Set–up and standing hypotheses}

\paragraph{Measurable spaces.}
Let $(\mathcal X ,\mathcal B_{\!X})$ and $(\mathcal Z ,\mathcal B_{\!Z})$ be measurable spaces. A fixed distribution $P_X$ on $(\mathcal X ,\mathcal B_{\!X})$ is given.

\paragraph{Correctness.}
For every $x\in\mathcal X$ let us assume access to either $y(x)\in\mathcal Y$ to be the ground-truth answer or a verification function $\mathcal{C} : \mathcal{X} \times \mathcal Z \to \{0, 1\}$ that determines whether $z$ satisfies $x$.
Define
\[
\corr(x,z)\;:=\;\one\!\bigl[\,\text{``the answer encoded in $z$ equals }y(x)\text{'' }\lor\text{``the answer encoded in $z$ is verified by }\mathcal C\text{ ''}\bigr]\in\{0,1\},
\qquad
(x,z)\in\mathcal X\times\mathcal Z.
\]

\paragraph{Policy classes.}
Throughout this section \emph{policies are deterministic measurable maps.}

\smallskip
\begin{itemize}
\item Helpful prover $\dH :\mathcal X\to\mathcal Z$.
\item Sneaky prover $\dS :\mathcal X\times\mathcal Z\to\mathcal Z$.
\item Verifier $V:\mathcal X\times\mathcal Z\to[0,1]$.
\end{itemize}

\paragraph{Per–input interaction.}
For every round, on one common probability space,
\[
x\sim P_X,\qquad
z_H=\dH(x),\qquad
z_S=\dS(x,z_H).
\]

\paragraph{Verifier loss.}
With $\sig(t)=\bigl(1+e^{-t}\bigr)^{-1}$ define
\[
\Delta V(x):=V(x,z_H)-V(x,z_S),\qquad
L_V(V,\dH,\dS)
:=-\EE_{x}\!\left[\log\sig\bigl(\Delta V(x)\bigr)\right].
\tag{1.1}
\]

\paragraph{Reward shape.}
A measurable function
$r:[0,1]\times\{0,1\}\to\mathbb R$
is fixed such that

\begin{enumerate}
\item[(R1)] (uniform alignment gap) there exists $\gamma>0$ with
      $r(s,1)\;-\;r(s',0)\;\ge\;\gamma\quad\forall\,s,s'\in[-B,B]$;
\item[(R2)] (strict monotonicity) for the map $s\mapsto r(s,1)$ there
      exists $L>0$ such that
      $r(s_1,1)-r(s_2,1)\;\le\;L\,|s_1-s_2|$ and $r(s_1,1)>r(s_2,1)$ whenever $s_1>s_2$.
\end{enumerate}

\paragraph{Utilities of the provers.}
\[
\begin{aligned}
U_H(V,\dH,\dS)
&:=\EE_x\!\bigl[r\bigl(V(x,\dH(x)),\corr(x,\dH(x))\bigr)\bigr],\\[2mm]
U_S(V,\dH,\dS)
&:=\EE_x\!\bigl[r\bigl(V(x,\dS(x,\dH(x))),1-\corr(x,\dS(x,\dH(x)))\bigr)\bigr].
\end{aligned}
\tag{1.2}
\]

\section*{2 \; Solution concepts}

\subsection*{2.1 \; Verifier-first $\varepsilon$–Stackelberg equilibrium}

\begin{definition}[ $(\varepsilon_V,\varepsilon_H,\varepsilon_S)$-VFSE]
Fix $(\varepsilon_V,\varepsilon_H,\varepsilon_S)\in\mathbb R_{\ge 0}^3$.
A triple $(V^\star,\dH^\star,\dS^\star)$ is an $(\varepsilon_V,\varepsilon_H,\varepsilon_S)$-verifier-first
Stackelberg equilibrium (VFSE) if

\begin{enumerate}
\item[(V)] (near-optimal leader)
      \[
      L_V\!\bigl(V^\star,\dH^\star,\dS^\star\bigr)
      \;\le\;
      \inf_{V'}L_V\!\bigl(V',\dH^{\operatorname{BR}}(V'),\dS^{\operatorname{BR}}(V',\dH^{\operatorname{BR}}(V'))\bigr)
      \;+\;\varepsilon_V,
      \]
      where the followers are the best responses below;
\item[(H)] (helpful $\varepsilon_H$–best response)
      \[
      U_H\!\bigl(V^\star,\dH^\star,\dS^\star\bigr)
      \;\ge\;
      \sup_{\dH'}U_H\!\bigl(V^\star,\dH',\dS^{\operatorname{BR}}(V^\star,\dH')\bigr)-\varepsilon_H;
      \]
\item[(S)] (sneaky per-input $\varepsilon_S$–best response)
      for $P_X$-a.e.\ $x$ one has
      \[
      U_S\!\bigl(V^\star,\dH^\star,
                 \dS^\star(\cdot\mid x,\dH^\star(x))\bigr)
      \;\ge\;
      \sup_{\mu\in\Delta(\mathcal Z)}
        U_S\!\bigl(V^\star,\dH^\star,\mu\bigr)-\varepsilon_S,
      \]
      where the supremum ranges over all measurable point-wise deviations
      $\mu(z)=\PP[z_S=z]$.\qedhere
\end{enumerate}
\end{definition}

\subsection*{2.2 \; $(\alpha,\beta)$–Sequential-Controllability}

\begin{definition}[Sequential-Controllability]
Given $(\alpha,\beta)\in[0,1]^2$
a triple $(V,\dH,\dS)$ is \emph{$(\alpha,\beta)$–Sequentially-Controllable} if

\[
\begin{aligned}
&\text{(C1) }
\EE_x\bigl[\corr(x,\dH(x))\bigr]\;\ge\;1-\alpha,\\[3pt]
&\text{(C2) }
\PP\!\bigl[\,V(x,\dH(x))>V\bigl(x,\dS(x,\dH(x))\bigr)\bigr]\;\ge\;1-\alpha,\\[3pt]
&\text{(S)}\quad
\forall\; \text{measurable } \dSt:\mathcal X\times\mathcal Z\to\mathcal Z, \\
&\quad
\PP\!\Bigl[\corr\!\bigl(x,\dSt(x,\dH(x))\bigr)=0
          \;\wedge\;
          V\bigl(x,\dH(x)\bigr)\le
          V\!\bigl(x,\dSt(x,\dH(x))\bigr)\Bigr]
          \;\le\;\beta.
\end{aligned}
\]
\end{definition}

\section*{3 \; Main theorem}

\begin{theorem}[Equivalence]
\label{thm:equivalence}
Assume (R1–R2) and $|V(x,z)|\le B$ for all $V\in\mathcal V$.
Define the \emph{moduli}
\[
\begin{aligned}
f_1(\varepsilon_V,\varepsilon_H,\varepsilon_S)
      &:=\Bigl( \alpha=\frac{\varepsilon_H}{\gamma}\;\vee\;\frac{\varepsilon_V}{\ln 2},
                 \;\beta=\frac{\varepsilon_S}{L\,B}\Bigr),\\[4pt]
f_2(\alpha,\beta)
      &:=\Bigl( \varepsilon_V=\alpha\ln 2,\;
                 \varepsilon_H=L\,B\,\alpha\;\vee\;\gamma\,\alpha,\;
                 \varepsilon_S=L\,B\,\beta\Bigr).
\end{aligned}
\]

\smallskip\noindent
\textbf{(i)}
If $(V^\star,\dH^\star,\dS^\star)$ is an
$(\varepsilon_V,\varepsilon_H,\varepsilon_S)$-VFSE,
then it is
$(\alpha,\beta)$–Sequentially-Controllable with $(\alpha,\beta)=
f_1(\varepsilon_V,\varepsilon_H,\varepsilon_S)$.

\smallskip\noindent
\textbf{(ii)}
Conversely, if $(V,\dH,\dS)$ is $(\alpha,\beta)$–Sequentially-Controllable and $\dS$ is a $\varepsilon_S'$–best response to $(V,\dH)$ point-wise, then $(V,\dH,\dS)$ is an $(\varepsilon_V,\varepsilon_H,\varepsilon_S)$-VFSE with $(\varepsilon_V,\varepsilon_H,\varepsilon_S) = f_2(\alpha,\beta)\;+\ (0,0,\varepsilon_S')$.
\\
Both statements are independent of any feature that $\dH$ and $\dS$ may share in their outputs.
\end{theorem}

\section*{4 \; Proof of Theorem~\ref{thm:equivalence}}

The proof is divided into the two implications.

\subsection*{4.1 \; From VFSE to Sequential-Controllability}

Let $(V^\star,\dH^\star,\dS^\star)$ satisfy Definition~2.1 with
tolerances $(\varepsilon_V,\varepsilon_H,\varepsilon_S)$.

\paragraph{Step 1: completeness (C1).}
Put
$p:=\PP\bigl[\corr(x,\dH^\star(x))=0\bigr]$.
Consider the alternative helper $\widehat{\dH}$ that, whenever possible,
outputs \emph{some} correct solution\footnote{If the set of correct
solutions is empty on a null set of $x$, we may modify $\widehat{\dH}$
arbitrarily there without affecting expectations; otherwise our
construction is undefined and both C1 and the conclusion are vacuous.}.
By (R1)

\[
U_H\!\bigl(V^\star,\widehat{\dH},\dS^{\star}\bigr)
-
U_H\!\bigl(V^\star,\dH^\star,\dS^{\star}\bigr)
\;\ge\;
\gamma\,p.
\]

Because $\dH^\star$ is an $\varepsilon_H$-best response,
$\gamma\,p\le\varepsilon_H$, hence
\[
p\le\frac{\varepsilon_H}{\gamma}\quad\Longrightarrow\quad
\EE[\corr(x,\dH^\star(x))]\ge 1-\frac{\varepsilon_H}{\gamma}.
\tag{4.1}
\]

\paragraph{Step 2: completeness (C2).}
Let
$q:=\PP\!\bigl[V^\star(x,\dH^\star(x))\le V^\star\bigl(x,\dS^\star(x,\dH^\star(x))\bigr)\bigr]$.
For every sample on which the inequality fails we have
$\sig(\Delta V)\le\frac12$ and thus $\log\sig(\Delta V)\le-\ln 2$.
On the remaining samples $\log\sig(\Delta V)\le 0$.
Consequently
\[
L_V\!\bigl(V^\star,\dH^\star,\dS^\star\bigr)
\;\ge\;
q\,\ln 2.
\]

If $q>\varepsilon_V/\ln 2$ then
$L_V(V^\star,\dH^\star,\dS^\star)>\varepsilon_V$,
yet the infimum of $L_V$ is $\le 0$
(achieved in the limit $\Delta V\to+\infty$),
contradicting verifier optimality (V).
Therefore
\[
q\;\le\;\frac{\varepsilon_V}{\ln 2},
\qquad
\PP\bigl[V^\star(x,\dH^\star(x))>V^\star(x,\dS^\star(x,\dH^\star(x)))\bigr]
\;\ge\;1-\frac{\varepsilon_V}{\ln 2}.
\tag{4.2}
\]

\paragraph{Step 3: soundness (S).}
Fix an arbitrary measurable $\dSt$.
Define the bad event
\[
E\;:=\;\Bigl\{
\corr\bigl(x,\dSt(x,\dH^\star(x))\bigr)=0
\ \wedge\
V^\star(x,\dH^\star(x))\le
V^\star\bigl(x,\dSt(x,\dH^\star(x))\bigr)
\Bigr\}.
\]

On $E$ the sneaky prover $\dSt$ secures a verifier score at least as
high as $\dS^\star$, and (by monotonicity (R2)) a strictly higher reward
whenever the inequality is strict.
Hence
\[
U_S\!\bigl(V^\star,\dH^\star,\dSt\bigr)
-
U_S\!\bigl(V^\star,\dH^\star,\dS^\star\bigr)
\;\ge\;
L\,B\,\PP[E].
\]

Because $\dS^\star$ is a point-wise $\varepsilon_S$–best response
condition~(S) in Definition~2.1 implies
$\PP[E]\le\varepsilon_S/(L\,B)$.

\paragraph{Conclusion.}
Combining \eqref{4.1}, \eqref{4.2} and the last inequality yields
Sequential-Controllability with
\[
\alpha
=\frac{\varepsilon_H}{\gamma}\ \vee\ \frac{\varepsilon_V}{\ln 2},
\qquad
\beta
=\frac{\varepsilon_S}{L\,B},
\]
i.e.\ $(\alpha,\beta)=f_1(\varepsilon_V,\varepsilon_H,\varepsilon_S)$.
\qed

\subsection*{4.2 \; From Sequential-Controllability to VFSE}

Assume now that $(V,\dH,\dS)$ is $(\alpha,\beta)$–controllable, and that
$\dS$ is, for every $(x,z_H)$, an $\varepsilon_S'$–best response.

\paragraph{Step 1: verifier near-optimality.}
On the good event of measure $1-\alpha$ we have
$V(x,\dH(x))>V(x,\dS(x,\dH(x)))$ and hence
$\sig(\Delta V)\ge \sig(0)=\tfrac12$,
so $\log\sig(\Delta V)\ge-\ln 2/2$.
On the bad event (measure $\alpha$) we
only know $\log\sig(\Delta V)\ge-\ln 2$.
Therefore
\[
L_V(V,\dH,\dS)
\;\le\;
\alpha\ln 2
\;=\;\varepsilon_V
\quad\text{with }\varepsilon_V:=\alpha\ln 2.
\]
Since the global infimum of $L_V$ is $\ge 0$,
$V$ is an $\varepsilon_V$-optimal leader.

\paragraph{Step 2: helpful near-best-response.}
Let $\dHt$ be any deviation.
Write
\[
\delta_c
:=\PP\bigl[\corr(x,\dHt(x))<\corr(x,\dH(x))\bigr]\le 1,
\qquad
\delta_s
:=\EE\bigl[V(x,\dHt(x))-V(x,\dH(x))\bigr].
\]
Reward difference, using (R1–R2):
\[
U_H(V,\dHt,\dS)-U_H(V,\dH,\dS)
\;\le\;
-\gamma\,\delta_c + L\,|\delta_s|.
\]
If $\dHt$ were to improve $U_H$ by more than
$\varepsilon_H:=\max\{\gamma\alpha,\,L\,B\,\alpha\}$,
then either $\delta_c>\alpha$ or $|\delta_s|>B\alpha$.
Both contradict (C1) or (C2) respectively (each fails on a set of
measure at most $\alpha$).
Hence $\dH$ is an $\varepsilon_H$-best response.

\paragraph{Step 3: sneaky near-best-response.}
For any measurable deviation $\dSt$,
\[
U_S(V,\dH,\dSt)-U_S(V,\dH,\dS)
\;\le\;
L\,B\,\PP[E]
\;\le\;
L\,B\,\beta,
\]
where $E$ is exactly the bad event in soundness.
Taking $\varepsilon_S:=L\,B\,\beta+\varepsilon_S'$ we obtain that
$\dS$ is an $\varepsilon_S$-best response.

\paragraph{Conclusion.}
Thus $(V,\dH,\dS)$ satisfies all three VFSE conditions with
\[
(\varepsilon_V,\varepsilon_H,\varepsilon_S)
=
f_2(\alpha,\beta)\;+\;(0,0,\varepsilon_S'),
\]
as claimed.  \qed

\section*{5 \; Remarks on common-prefix collusion}

The verifier loss~\eqref{1.1} depends only on the \emph{difference}
$V(x,z_H)-V(x,z_S)$.
If $\dH$ and $\dS$ collude by prepending an identical string
$P_c$ to $z_H$ and $z_S$, the value of $L_V$ is unchanged.
Hence any feature shared by both provers is irrelevant for optimisation;
the formal inequalities above hold \emph{regardless} of such collusion.

\section*{6 \; Extension to stochastic neural policies}

Replacing deterministic maps by neural networks that sample from
$\Delta(\mathcal Z)$ simply inserts an outer expectation
over the networks’ internal randomness into every definition.
All bounds derived in Section~4 are linear in probabilities
and Lipschitz in verifier scores; they therefore continue to hold
verbatim after the additional expectation.
Existence of approximate equilibria then follows from standard
compactness arguments for finite-parameter networks or,
operationally, from convergence to approximate first-order stationarity
in stochastic gradient training.

\bigskip
\begin{flushright}
\qedsymbol
\end{flushright}
